\section{Discussion}
\label{sec:discussion}

\paragraph{Is this just instruction following?}
In part it is: the model follows the schema over ``List all'', and Section~\ref{sec:scoring} gives our reasons for counting that as the failure.
The controls show that the minimum does this, not the schema or any number in it.
Under a shown maximum 35.7\% of outputs over-count, against 34.5\% with no count and 55.4\% under a shown minimum.
A minimum adds no entries when the document already holds enough items.
Told that a shorter list is fine, seven of eight GPU models shorten their lists against the plain minimum, so most can stop.
What goes wrong is that the model meets the count with whatever text is at hand instead of saying the list is short.

\paragraph{Why it is easy to miss.}
Most of the padding is taken from the document: it parses, meets the schema, and about half of it passes a source check (Section~\ref{sec:what}), so the error shows in full only against an answer list.
This matches what \citet{han2023empirical} and \citet{qian2026unreliable} find for extraction without a minimum count, and differs from PhantomFill, where the inputs hold nothing to copy \citep{usman2026phantomfill}.

\begin{table}[t]
  \centering
  \small
  \fbox{\begin{minipage}{0.92\columnwidth}
  \textbf{For schema authors}:
  \begin{itemize}
    \setlength\itemsep{0pt}
    \item Leave \texttt{minItems} off lists whose length the document decides (over-count 34.5\% $\to$ 55.4\% shown, 90.1\% enforced).
    \item If recall matters most, an enforced minimum misses the fewest true items but pads most; prose ``up to $k$'' balances the two best (F1 0.60\pre{} vs 0.53\pre{}).
    \item Do not rely on a source check: about half of the padding passes it.
    \item If a minimum must stay, add ``fewer is fine'' and test it: some models still lengthen their lists, and recall falls.
    \item An ``N/A'' option under enforcement helps but is not a fix (90.1\% $\to$ 51.0\% padded).
    \item Do not re-ask with the bare validation error (Qwen3.5-9B: 10.0\% $\to$ 69.0\% over-count).
    \item \texttt{maxItems} barely changes over-count but shortens lists and leaves out more true items.
  \end{itemize}
  \end{minipage}}
  \caption{Guidance from the open-weight models tested.}
  \label{tab:guidance}
\end{table}

\section{Conclusion}
\label{sec:conclusion}

When a document holds fewer items than a schema's minimum, most of the open-weight models tested treat the count as a target and pad their lists: a shown minimum raises the share of over-counting lists from 34.5\% to 55.4\%, and an ``N/A'' option under an enforced minimum still leaves half of them padded.
The added entries are mostly real text taken from the document, so about half of them pass a source check.
A ``fewer is fine'' line undoes the effect on average at some cost in recall, while re-asking with the validation error makes padding far more common, and thinking does not help.
